
\documentclass[a4paper,12pt]{article}
\usepackage[italian]{babel}
\usepackage{listings}
\usepackage{geometry}
\geometry{margin=2.5cm}

\title{Verifica Conclusica del Corso di Basi di Dati}
\author{Prof. Fedeli Massimo  - Fabbrica Digitale 4.0}
\date{}

\begin{document}
\maketitle

Si consideri il seguente schema logico relativo ad un sistema di bike sharing cittadino.

CLIENTE(\underline{idCliente}, nome, cognome, email, citta)

BICI(\underline{idBici}, modello, tipo, stato, idStazione)

STAZIONE(\underline{idStazione}, nome, citta, numeroPosti)

NOLEGGIO(\underline{idNoleggio}, dataInizio, dataFine, idCliente, idBici)

MANUTENZIONE(\underline{idManutenzione}, dataManutenzione, tipoIntervento, idBici)

\section{Istruzione}
\begin{itemize}
	\item Scrivere le queries all'interno di un file di testo denominato nel seguente modo:  cognome-nome.md
	\item Inviare per email al docente il file.
\end{itemize}

\section*{Esercizio 1 - Query SQL}

Scrivere le seguenti interrogazioni SQL.

\begin{enumerate}

\item Elencare nome e cognome dei clienti ordinati alfabeticamente per cognome e, a parità di cognome, per nome.

\item Elencare idBici, modello e tipo delle biciclette elettriche ordinate prima per modello e poi per idBici decrescente.

\item Elencare nome e numeroPosti delle stazioni situate nella città di Milano ordinate per numeroPosti decrescente.

\item Elencare modello, tipo e nome della stazione in cui si trova ogni bicicletta.

\item Calcolare quante biciclette sono presenti in ogni stazione mostrando idStazione e numero di biciclette.

\item Elencare idNoleggio, dataInizio, nome e cognome del cliente che ha effettuato il noleggio.

\item Elencare modello della bicicletta, nome del cliente e dataInizio per ogni noleggio effettuato.

\item Calcolare per ogni cliente il numero totale di noleggi effettuati mostrando nome, cognome e numero di noleggi.

\item Elencare i clienti che non hanno mai effettuato alcun noleggio.

\item Elencare idBici e modello delle biciclette che non sono mai state noleggiate.

\item Calcolare per ogni bicicletta il numero totale di noleggi effettuati mostrando idBici e numero di noleggi.

\item Trovare le stazioni che possiedono più di 3 biciclette mostrando nome della stazione e numero di biciclette.

\item Trovare i clienti che hanno effettuato almeno 2 noleggi mostrando nome, cognome e numero di noleggi.

\item Elencare i modelli delle biciclette che sono state noleggiate almeno una volta ma che non hanno mai subito interventi di manutenzione.

\item Trovare l'id della bicicletta che è stata noleggiata il maggior numero di volte insieme al numero di noleggi.
\end{enumerate}

\section*{Esercizio 2 - SQL DDL}

Scrivere le istruzioni SQL per creare le seguenti tabelle specificando chiavi primarie e chiavi esterne:

\begin{itemize}
\item CLIENTE
\item NOLEGGIO (con chiave esterna verso CLIENTE)
\end{itemize}

\section*{Esercizio 3 - Viste}

Scrivere una vista denominata \texttt{vista\_noleggi\_clienti} che mostri:

\begin{itemize}
\item nome cliente
\item cognome cliente
\item modello bici
\item data inizio noleggio
\end{itemize}

\end{document}
